\chapter*{\centering\Large{TÓM TẮT ĐỒ ÁN}}
\addcontentsline{toc}{chapter}{Tóm tắt đồ án}

Trong bối cảnh các cuộc tấn công mạng nhắm vào ứng dụng Web ngày càng phức tạp và gia tăng nhanh chóng, việc kiểm thử xâm nhập đóng vai trò then chốt trong việc bảo vệ an toàn thông tin. Tuy nhiên, quy trình kiểm thử thủ công đòi hỏi nhiều nguồn lực, thời gian và chuyên môn cao, trong khi các công cụ tự động truyền thống thường thiếu khả năng suy luận linh hoạt và không hiểu được ngữ cảnh nghiệp vụ, dẫn đến tỷ lệ bỏ sót lỗ hổng cao. Đồ án này trình bày việc thiết kế và hiện thực một hệ thống kiểm thử xâm nhập ứng dụng Web tự động, ứng dụng kiến trúc đa tác tử AI kết hợp sức mạnh của mô hình ngôn ngữ lớn Gemini nhằm giải quyết các hạn chế nêu trên.

Hệ thống được thiết kế tuân thủ nghiêm ngặt bộ tiêu chuẩn OWASP Web Security Testing Guide, WSTG, phiên bản 4.2, bao phủ 105 kịch bản kiểm thử trải đều trên 12 nhóm lỗ hổng bảo mật. Để đảm bảo tính chính xác và chống hiện tượng "ảo giác" thường gặp ở mô hình ngôn ngữ lớn, hệ thống áp dụng bốn trụ cột công nghệ cốt lõi:, 1, Cơ chế RAG truy vấn tri thức chuẩn từ OWASP theo phương pháp tìm kiếm tương đồng vector;, 2, Đồ thị tri thức động đóng vai trò như bộ não lưu trữ, liên kết và tái sử dụng dữ liệu trinh sát, như cổng mạng, cấu trúc thư mục, loại công nghệ, xuyên suốt giữa các bài kiểm thử;, 3, Kiến trúc đa tác tử phân quyền chặt chẽ giữa tác tử chỉ huy, Planner Agent,, tác tử trinh sát, Recon Agent,, tác tử khai thác, Exploit Agent, và tác tử xác thực, Verifier,, kết hợp với cơ chế xác thực tự động nhằm giảm thiểu tỷ lệ cảnh báo sai;, 4, Giao thức MCP đóng vai trò cầu nối an toàn giữa AI và các công cụ kiểm thử thực tế, Nmap, SQLMap, ZAP, v.v.,, đảm bảo tính module và khả năng mở rộng.

Kết quả thực nghiệm trên OWASP Juice Shop - một ứng dụng Web mô phỏng có chủ đích chứa nhiều lỗ hổng - cho thấy hệ thống có khả năng tự động hoàn thành 105 bài kiểm thử, đồng thời phát hiện thành công các lỗ hổng nghiêm trọng như SQL Injection, Cross-Site Scripting, XSS,, và lộ thông tin cấu hình nhạy cảm. Hệ thống cũng thể hiện hiệu ứng tích lũy tri thức rõ rệt khi dữ liệu từ các bước thăm dò trước được sử dụng làm đầu vào chính xác cho các bước tấn công sau, qua đó mô phỏng thành công luồng tư duy của một chuyên gia kiểm thử bảo mật thực thụ. Hướng nghiên cứu này mở ra tiềm năng lớn trong việc tích hợp AI vào quy trình DevSecOps để liên tục phát hiện và khắc phục lỗ hổng trong vòng đời phát triển phần mềm.